\begin{circuitikz}[scale=0.95, transform shape]
  % Referencia de masa
  \draw (0,0) node[ground] {}
    to[sV, l=$v_s$] (0,3)
    to[R, l=$R_s$] (1.8,3)
    to[C, l=$C_1$, -*] (3.5,3) coordinate (base_node);

  % Divisor resistivo de entrada
  \draw (3.5,6.5) node[vcc] {$V_{CC}$} coordinate (vcc_node)
    to[R, l=$R_2$, -*] (base_node)
    to[R, l=$R_1$, *-] (3.5,0) node[ground] {};

  % Transistor BJT NPN
  \draw (base_node) to[short] (4.2,3) node[npn, anchor=B] (Q1) {};
  \node[right=4pt] at (Q1.right) {$Q_1$ (BC546B)};

  % Malla de colector
  \draw (Q1.C) to[short, -*] ++(0,0.25) coordinate (coll_node)
    to[R, l_=$R_C$] ++(0,2) coordinate (aux) to [short] (vcc_node |- aux);

  % Malla de emisor con capacitor de desacoplo
  \draw (Q1.E) to[short, -*] ++(0,-0.5) coordinate (emit_node);
  \draw (emit_node) to[R, l=$R_E$] (5,0) node[ground] {};
\draw (emit_node) to[short] ++(1.5,0) coordinate(aux)
      to[C, l=$C_E$] (aux |- 0,0) node[ground] {};

  % Red de acoplo de salida y carga
  \draw (coll_node) to[C, l=$C_2$] ++(4,0) coordinate (out_node)
    to[short, -o] ++(1,0) node[above] {$v_L$};
  \draw (out_node) to[R, l=$R_L$, *-] (9,0) node[ground] {};
\end{circuitikz}
